%\chapter*{Conclusion}
%\addcontentsline{toc}{chapter}{Conclusion}
%
%\paragraph{Introduction de la conclusion.}
%Cette thèse s’est concentrée sur la caractérisation expérimentale et théorique des gaz de bosons 1D, systèmes quantiques intégrables où l’état relaxé n’est pas thermique mais défini par la distribution de rapidités. L’objectif principal était de sonder cette distribution et de comparer les observations avec les prédictions de la théorie \gls{GHD}.
%
%\paragraph{Résumé des travaux réalisés.}
%Au cours de ces trois années, mon travail s’est articulé autour de plusieurs axes complémentaires. Sur le plan théorique, j’ai approfondi ma compréhension du modèle de \gls{LL} et de l’\gls{BA}, ainsi que du formalisme des systèmes intégrables, en particulier la \gls{GGE} et la \gls{GHD}. J’ai particulièrement apprécié l’étude du formalisme inverse de Scattering, de l’ansatz algébrique de Bethe, ainsi que des liens avec la théorie des matrices aléatoires présentés dans ce mémoire. Ces travaux théoriques ont permis de mieux comprendre la dynamique des gaz de bosons 1D et d’illustrer concrètement mes cours dans le cadre de ma thèse.
%
%Sur le plan numérique et expérimental, j’ai réalisé des simulations GHD de protocoles de déformation des bords et de calculs de fluctuations, sans mesurer directement les fluctuations expérimentales. J’ai également participé à l’initialisation et à la mise en place d’un piégeage dipolaire, première étape nécessaire pour réaliser des mesures locales sur le nuage atomique
%
%\paragraph{Apports principaux.}
%Cette thèse a permis de développer et de mettre en œuvre une approche combinant théorie, simulations numériques et préparation expérimentale pour étudier les gaz de bosons 1D. Mon apport principal réside dans la compréhension approfondie des systèmes intégrables, l’implémentation de simulations GHD pour des protocoles dynamiques, ainsi que dans la mise en place initiale d’un piégeage dipolaire ouvrant la voie à des mesures locales futures.
%
%\paragraph{Perspectives.}
%Les travaux futurs pourront se concentrer sur l’optimisation du piégeage dipolaire pour permettre des mesures expérimentales directes des fluctuations de rapidités. L’étude des effets de pertes atomiques et de leur influence sur les distributions de rapidité constitue également une piste prometteuse. Ces développements permettront de mieux explorer les dynamiques hors équilibre et d’approfondir la compréhension des systèmes intégrables dans des conditions expérimentales contrôlées.


%\chapter*{Conclusion}
%\addcontentsline{toc}{chapter}{Conclusion}
%
%\section*{Introduction de la conclusion}
%Dans cette thèse, j’ai exploré les gaz de bosons 1D intégrables, où l’état relaxé ne suit pas la statistique Gibbs classique, mais est caractérisé par la distribution de rapidités. L’objectif central était de sonder cette distribution spatialement résolue et de confronter les données expérimentales aux prédictions de la théorie d’Hydrodynamique Généralisée (\gls{GHD}).
%
%\section*{Résumé des travaux réalisés}
%
%Cette thèse a permis de combiner théorie, simulation et expérimentation pour mieux comprendre et manipuler les gaz de bosons 1D intégrables. Mes contributions se situent dans trois axes :
%
%\subsection*{Approche théorique}
%J’ai approfondi mes compréhension dans les fondements du modèle de \acrlong{LL} et de l’Ansatz de Bethe (\gls{BA}), ainsi que leur encadrement par le cadre théorique des systèmes intégrables : \gls{GGE}, \gls{TBA} et \gls{GHD}. Ces mécanismes, au carrefour de la théorie quantique et de l’hydrodynamique quantique, ont renforcé ma compréhension conceptuelle, tout en illustrant mes cours dans un contexte applicatif.
%
%%\subsection*{Simulations numériques et modélisation}
%%J’ai mis en place des simulations \gls{GHD} pour des protocoles dynamiques, notamment des déformations de profil de densité débutant en écart de bord, et ai entamé une exploration numérique des fluctuations de distribution de rapidité, via le principe de fluctuation--réponse et un échantillonnage Monte Carlo --- sans toutefois disposer de mesures expérimentales directes sur ces fluctuations pendant cette thèse.
%
%\subsection*{Simulations numériques et modélisation}
%Dans cette thèse, j’ai développé des simulations \gls{GHD} pour étudier la dynamique des gaz de bosons 1D, notamment dans le cadre de protocoles impliquant des déformations locales des profils de densité. En parallèle, j’ai calculé les fluctuations de la distribution de rapidité en utilisant le formalisme \gls{TBA}, et validé ces résultats par des vérifications numériques fondées sur le principe de fluctuation–réponse.  
%
%Afin d’explorer la faisabilité de l’échantillonnage du \gls{GGE}, j’ai mis en œuvre des algorithmes Monte Carlo adaptés. Ces simulations permettent non seulement d’étudier les corrélations mais aussi permettra d'étudier des  fluctuations d’ordre supérieur. Elles constituent un cadre théorique robuste pour anticiper et comparer avec de futures mesures expérimentales locales, bien que celles-ci n’aient pas été réalisées au cours de cette thèse.
%
%Afin d’explorer la faisabilité de l’échantillonnage du \gls{GGE}, j’ai amorcé des simulations Monte Carlo adaptées. Ces outils fournissent un cadre théorique robuste pour anticiper de futures mesures expérimentales locales, bien que celles-ci n’aient pas été réalisées au cours de cette thèse.
%
%
%\subsection*{Mise en place expérimentale}
%Sur le plan expérimental, j’ai initié la mise en place d’un piégeage dipolaire en complément de la puce atomique, visant à préparer des configurations initiales de gaz atomiques peu conventionnelles. Cette étape permet de générer des potentiels longitudinaux complexes et modulables, ouvrant la voie à l’étude future de distributions de rapidité spatialement résolues et de dynamiques hors équilibre dans des états initialement « exotiques ».
%
%
%%\section*{Apports principaux}
%%Cette thèse a permis de combiner théorie, simulation et expérimentation pour mieux comprendre et manipuler les gaz de bosons 1D intégrables. Mes contributions se situent dans trois axes :
%%\begin{itemize}[label = $\bullet$]
%%  \item \textbf{Approfondissement théorique} : maîtrise du formalisme des gaz 1D, \gls{GGE}, \gls{TBA} et \gls{GHD}.
%%  \item \textbf{Simulation GHD avancée} : description dynamique de l’évolution de la distribution de rapidités, y compris en régime transitionnel.
%%  \item \textbf{Préparation expérimentale du confinement dipolaire} : étape indispensable pour des mesures futures de distributions locales.
%%\end{itemize}
%
%%\section*{Perspectives et nouvelles pistes à explorer}
%%Voici des axes enrichis et affinés pour prolonger et valoriser ces travaux :
%%\begin{itemize}[label = $\bullet$]
%%  \item \textbf{Optimisation du confinement dipolaire} : calibrer précisément la géométrie, la puissance et la stabilité du piège pour accéder à des régimes où les mesures locales de fluctuations deviennent possibles.
%%  \item \textbf{Simulation Monte Carlo de fluctuations d’ordre supérieur} : développer des algorithmes plus efficaces pour générer \gls{GGE} et analyser les corrélations de fluctuations de haut ordre.
%%  \item \textbf{Inclusion des effets de pertes atomiques} : s’appuyer sur la récente modélisation de l’évolution de la distribution de rapidités sous pertes atomiques pour étudier les déviations non thermiques et les intégrer dans le cadre GHD \cite{10.21468/SciPostPhys.9.4.044}.
%%  \item \textbf{Étude des pertes multibody en présence de confinement} : poursuivre des simulations numériques adaptées à la géométrie expérimentale continue afin d’évaluer la robustesse des distributions de rapidité.
%%  %\item \textbf{Thermométrie par fluctuations localisées} : appliquer des méthodes inspirées de la microscopie quantique (quantum gas microscopy) afin de mesurer température ou distribution de rapidités sans hypothèse de globalité.
%%  \item \textbf{Évaluation expérimentale des fluctuations de rapidité} : envisager des mesures directes des fluctuations en densité après expansion ou \emph{in situ} et les comparer aux prédictions théoriques et Monte Carlo.
%%\end{itemize}
%%
%%\section*{Faisabilité et réflexion critique}
%%\begin{itemize}[label = $\bullet$]
%%  \item \textbf{Simulation Monte Carlo} : réalisable avec l’architecture déjà développée, mais coûteuse en temps de calcul pour un échantillonnage fin et stable.
%%  %\item \textbf{Pertes atomiques et dynamique GHD} : modélisable grâce aux résultats analytiques existants et aux codes disponibles.
%%  \item \textbf{Contrôle expérimental du piégeage dipolaire} : accessible mais exigeant en termes d’optique, d’alignement et de stabilité.
%%  \item \textbf{Mesure de fluctuations} : techniquement complexe (résolution spatiale et statistique élevée), mais rendue envisageable avec les techniques modernes d’imagerie à haute résolution.
%%  %\item \textbf{Thermométrie par fluctuations} : prometteuse pour sonder les états hors équilibre, mais nécessitant une adaptation aux gaz de bosons 1D et une validation expérimentale robuste.
%%\end{itemize}
%
%
%\section*{Perspectives, faisabilité et réflexion critique}
%
%%\paragraph{Optimisation du confinement dipolaire.} 
%%Un calibrage précis de la géométrie, de la puissance et de la stabilité du piège permettrait d’accéder à des régimes où les mesures locales de fluctuations deviennent possibles. Cette perspective est accessible expérimentalement, mais exigeante en termes d’optique, d’alignement et de stabilité.
%
%\paragraph{Développement du piégeage dipolaire.} 
%La poursuite de la réalisation d’un piège dipolaire plus abouti offrirait une plus grande liberté dans la préparation des systèmes, tout en garantissant une meilleure reproductibilité des conditions expérimentales. Une telle maîtrise constitue un prérequis essentiel, par exemple pour envisager la mesure de fluctuations, mais elle demeure exigeante en termes d’optique, d’alignement et de stabilité.
%
%
%%\paragraph{Simulation Monte Carlo de fluctuations d’ordre supérieur.} 
%%Le développement d’algorithmes plus efficaces pour générer des états \gls{GGE} et analyser des corrélations de fluctuations de haut ordre constitue une piste prometteuse. %Cette approche est réalisable avec l’architecture déjà développée, mais reste coûteuse en temps de calcul pour obtenir un échantillonnage fin et stable.
%%
%%\paragraph{Inclusion des effets de pertes atomiques.} 
%%La modélisation récente de l’évolution de la distribution de rapidités sous pertes atomiques ouvre la voie à l’étude de déviations non thermiques, intégrables dans le cadre de la GHD \cite{10.21468/SciPostPhys.9.4.044}. Ces analyses demanderaient toutefois une validation numérique et expérimentale poussée.
%%
%%\paragraph{Étude des pertes multibody en présence de confinement.} 
%%La poursuite de simulations numériques adaptées à la géométrie expérimentale continue permettrait d’évaluer la robustesse des distributions de rapidité vis-à-vis des processus de pertes. Ce projet reste exigeant, mais techniquement envisageable.
%%
%%
%%\paragraph{Optimisation des simulations Monte Carlo pour les fluctuations d’ordre supérieur et les pertes.} 
%%Une piste prometteuse consiste à améliorer les méthodes de simulation Monte Carlo afin de générer plus efficacement des états \gls{GGE} et d’accéder à des corrélations de fluctuations d’ordre supérieur. Dans ce cadre, l’inclusion des effets de pertes atomiques, qu’il s’agisse de pertes simples ou multibody en présence de confinement, constitue un enjeu essentiel : ces processus dissipatifs, bien que distincts des fluctuations proprement dites, influencent significativement la dynamique des distributions de rapidité et pourraient révéler de nouvelles déviations non thermiques. La mise en œuvre de telles simulations numériques, adaptées à la géométrie expérimentale continue, reste exigeante mais offrirait un outil puissant pour confronter théorie et expérience.
%%
%%
%%\paragraph{Évaluation expérimentale des fluctuations de rapidité.} 
%%Des mesures directes des fluctuations en densité, après expansion ou \emph{in situ}, constitueraient une avancée majeure. Bien que techniquement complexes — nécessitant une résolution spatiale et statistique élevée — elles deviennent aujourd’hui envisageables grâce aux techniques modernes d’imagerie à haute résolution, et offriraient une comparaison directe avec les prédictions théoriques et Monte Carlo.
%
%
%\paragraph{Optimisation des simulations et étude des pertes.} 
%Une piste prometteuse consiste à améliorer les méthodes de simulation Monte Carlo afin de générer plus efficacement des états \gls{GGE} et d’accéder à des corrélations de fluctuations d’ordre supérieur. Dans ce cadre, l’inclusion des effets de pertes atomiques constitue un enjeu central : la modélisation récente de l’évolution de la distribution de rapidités sous pertes \cite{10.21468/SciPostPhys.9.4.044} ouvre en effet la voie à l’étude de déviations non thermiques, intégrables dans le cadre de la GHD. La prise en compte des pertes multibody, notamment dans une géométrie confinée, permettrait en outre d’évaluer la robustesse des distributions de rapidité face à ces processus dissipatifs. Enfin, la confrontation de ces résultats théoriques et numériques à des données expérimentales demeure une perspective majeure : des mesures directes des fluctuations de densité, après expansion ou \emph{in situ}, bien que techniquement complexes et exigeant une résolution spatiale et statistique élevée, deviennent aujourd’hui envisageables grâce aux techniques modernes d’imagerie à haute résolution, et offriraient une comparaison directe avec les prédictions issues des approches Monte Carlo.

\chapter*{Conclusion}
\addcontentsline{toc}{chapter}{Conclusion}
%\minitoc

\section*{Introduction de la conclusion}
Dans cette thèse, j’ai exploré les gaz de bosons 1D intégrables, où l’état relaxé ne suit pas la statistique Gibbs classique, mais est caractérisé par la distribution de rapidités. L’objectif central était de sonder cette distribution spatialement résolue et de confronter les données expérimentales aux prédictions de la théorie d’Hydrodynamique Généralisée (\gls{GHD}).

\section*{Résumé des travaux réalisés}

Cette thèse a permis de combiner théorie, simulation et expérimentation pour mieux comprendre et manipuler les gaz de bosons 1D intégrables. Mes contributions se situent dans trois axes :

\subsection*{Approche théorique}
J’ai approfondi mes compréhensions dans les fondements du modèle de \acrlong{LL} et de l’Ansatz de Bethe (\gls{BA}), ainsi que leur encadrement par le cadre théorique des systèmes intégrables : \gls{GGE}, \gls{TBA} et \gls{GHD}. Ces mécanismes, au carrefour de la théorie quantique et de l’hydrodynamique quantique, ont renforcé ma compréhension conceptuelle, tout en illustrant mes cours dans un contexte applicatif.


\subsection*{Simulations numériques et modélisation}
Dans cette thèse, j’ai développé des simulations \gls{GHD} pour étudier la dynamique des gaz de bosons 1D, notamment dans le cadre de protocoles impliquant des déformations locales des profils de densité. En parallèle, j’ai calculé les fluctuations de la distribution de rapidité en utilisant le formalisme \gls{TBA}, et validé ces résultats par des vérifications numériques fondées sur le principe de fluctuation–réponse.  

Afin d’explorer la faisabilité de l’échantillonnage du \gls{GGE}, j’ai amorcé des simulations Monte Carlo adaptées. Ces outils fournissent un cadre théorique robuste pour anticiper de futures mesures expérimentales locales, bien que celles-ci n’aient pas été réalisées au cours de cette thèse.


\subsection*{Mise en place expérimentale}
Sur le plan expérimental, j’ai initié la mise en place d’un piégeage dipolaire en complément de la puce atomique, visant à préparer des configurations initiales de gaz atomiques peu conventionnelles. Cette étape permet de générer des potentiels longitudinaux complexes et modulables, ouvrant la voie à l’étude future de distributions de rapidité spatialement résolues et de dynamiques hors équilibre dans des états initialement « exotiques ».


\section*{Perspectives, faisabilité et réflexion critique}


\paragraph*{Développement du piégeage dipolaire.}
La poursuite de la réalisation d’un piège dipolaire plus abouti offrirait une plus grande liberté dans la préparation des systèmes, tout en garantissant une meilleure reproductibilité des conditions expérimentales. Une telle maîtrise constitue un prérequis essentiel, par exemple pour envisager la mesure de fluctuations, mais elle demeure exigeante en termes d’optique, d’alignement et de stabilité.


\paragraph*{Optimisation des simulations et étude des pertes.}
Une piste prometteuse consiste à améliorer les méthodes de simulation Monte Carlo afin de générer plus efficacement des états \gls{GGE} et d’accéder à des corrélations de fluctuations d’ordre supérieur. Dans ce cadre, l’inclusion des effets de pertes atomiques et la modélisation récente de l’évolution de la distribution de rapidités sous pertes \cite{10.21468/SciPostPhys.9.4.044} constitue un enjeu central : ouvrir la voie à l’étude de déviations non thermiques, intégrables dans le cadre de la GHD. La prise en compte des pertes multibody, notamment dans une géométrie confinée, permettrait en outre d’évaluer la robustesse des distributions de rapidité face à ces processus dissipatifs. La confrontation de ces résultats, théoriques et numériques,  à des données expérimentales demeure une perspective majeure.

Bien que techniquement complexes et exigeant une résolution spatiale et statistique élevée, des mesures directes des fluctuations de densité, après expansion ou \emph{in situ},deviennent aujourd’hui envisageables grâce aux techniques modernes d’imagerie à haute résolution, et offriraient une comparaison directe avec les prédictions issues des approches Monte Carlo.
